\documentclass[10pt,journal,compsoc]{IEEEtran}
\IEEEoverridecommandlockouts

% -------------------- Essential Packages --------------------
\usepackage[cmex10]{amsmath}
\usepackage{amssymb,amsfonts,amsthm,mathtools,bm}
\usepackage{algorithm}
\usepackage{algorithmic}
\usepackage{graphicx}
\usepackage{textcomp}
\usepackage{xcolor}
\ifCLASSOPTIONcompsoc
  \usepackage[nocompress]{cite}
\else
  \usepackage{cite}
\fi
\usepackage{url}
\usepackage{booktabs}
\usepackage{multirow}
\usepackage{array}
\usepackage{tikz}
\usepackage{pgfplots}
\pgfplotsset{compat=1.17}
\usetikzlibrary{positioning,arrows,shapes,calc}
\usepackage{soul}
\usepackage{enumitem}
\usepackage{listings}
\usepackage{microtype}
\usepackage[caption=false,font=footnotesize]{subfig}
\usepackage[hidelinks]{hyperref}

\usepackage[nameinlink,capitalise,noabbrev]{cleveref}

\usepackage{tabularx}

% -------------------- Custom Commands --------------------
\newcommand{\R}{\mathbb{R}}
\newcommand{\E}{\mathbb{E}}
\newcommand{\N}{\mathcal{N}}
\newcommand{\F}{\mathcal{F}}
\newcommand{\G}{\mathcal{G}}
\newcommand{\D}{\mathcal{D}}
\newcommand{\A}{\mathcal{A}}
\newcommand{\X}{\mathcal{X}}
\newcommand{\Y}{\mathcal{Y}}
\newcommand{\Z}{\mathcal{Z}}
\newcommand{\s}{\mathcal{S}}
\newcommand{\Tcal}{\mathcal{T}}
\newcommand{\Kcal}{\mathcal{K}}
\newcommand{\Mcal}{\mathcal{M}}
\newcommand{\Q}{\mathcal{Q}}
\newcommand{\U}{\mathcal{U}}
\newcommand{\V}{\mathcal{V}}
\newcommand{\W}{\mathcal{W}}
\newcommand{\h}{\mathcal{H}}
\newcommand{\softplus}{\operatorname{softplus}}
\newcommand{\Softmax}{\operatorname{Softmax}}

\DeclareMathOperator*{\argmin}{arg\,min}
\DeclareMathOperator*{\argmax}{arg\,max}
\DeclareMathOperator{\Tr}{Tr}
\DeclareMathOperator{\diag}{diag}
\DeclareMathOperator{\rank}{rank}
\DeclareMathOperator{\vecop}{vec}
\DeclareMathOperator{\prox}{prox}
\DeclareMathOperator{\KL}{KL}
\DeclareMathOperator{\ELBO}{ELBO}

% -------------------- Theorem Environments --------------------
\newtheorem{theorem}{Theorem}
\newtheorem{lemma}[theorem]{Lemma}
\newtheorem{proposition}[theorem]{Proposition}
\newtheorem{corollary}[theorem]{Corollary}
\newtheorem{definition}{Definition}
\newtheorem{assumption}{Assumption}
\newtheorem{remark}{Remark}

% Tighten spacing around floats
\setlength{\textfloatsep}{5pt}% space between text and floats
\setlength{\floatsep}{5pt}% space between two floats
\setlength{\intextsep}{5pt}% space above and below in-text floats
\setlength{\skip\footins}{8pt}% space above footnotes

% -------------------- Title & Authors --------------------
\title{SDE-TGNN: Stochastic Differential Equation Temporal Graph Neural Network for Multi-Domain Network Intrusion Detection with Principled Uncertainty Quantification}

\author{%
Roger~Nick~Anaedevha,~\IEEEmembership{Student Member,~IEEE}\IEEEauthorrefmark{1},
Alexander~Gennadevich~Trofimov\IEEEauthorrefmark{1},
and~Yuri~Vladimirovich~Borodachev\IEEEauthorrefmark{2}\\
\IEEEauthorrefmark{1}National Research Nuclear University MEPhI (Moscow Engineering Physics Institute), Moscow 115409, Russia\\
\IEEEauthorrefmark{2}Artificial Intelligence Research Center, National Research Nuclear University MEPhI, Moscow 115409, Russia\\
Corresponding author: \href{mailto:ar006@campus.mephi.ru}{ar006@campus.mephi.ru}
}

% -------------------- Document --------------------
\begin{document}
\maketitle

\begin{abstract}
Network intrusion detection systems (NIDS) face critical challenges in handling heterogeneous network domains, quantifying prediction uncertainty, and providing robustness guarantees against adversarial attacks. We propose SDE-TGNN, a novel deep learning framework that combines Temporal Graph Neural Networks with Stochastic Differential Equations (SDEs) and Fokker-Planck-based analytical uncertainty propagation. Our model learns continuous-time network state evolution through drift and diffusion functions, capturing both deterministic trends and stochastic variability in traffic patterns. Unlike Monte Carlo sampling approaches, we propagate the first two moments of the state distribution analytically via the Fokker-Planck equation, achieving efficient $O(d^2)$ uncertainty quantification. We integrate SDE trajectories at multiple temporal scales to capture both short-term anomalies and long-term attack campaigns. Evaluated on six heterogeneous datasets spanning cloud, industrial IoT, container, enterprise, IoT, and hybrid network domains (totaling 53+ million samples), SDE-TGNN achieves state-of-the-art detection performance with a weighted average F1 score of 98.7\%, outperforming eight baseline models including classical ML, deep learning, graph neural networks, continuous-depth models, and Bayesian approximations. The model provides calibrated uncertainty estimates (ECE $<$ 0.05) and certified adversarial robustness guarantees via randomized smoothing. Our unified architecture demonstrates strong cross-domain generalization without domain-specific modifications, making it practical for real-world heterogeneous network security deployments.
\end{abstract}

\begin{IEEEkeywords}
Network Intrusion Detection, Stochastic Differential Equations, Temporal Graph Neural Networks, Fokker-Planck Equation, Uncertainty Quantification, Adversarial Robustness, Multi-Domain Learning
\end{IEEEkeywords}

\section{Introduction}
\IEEEPARstart{N}{etwork} intrusion detection systems constitute a critical defense layer in modern cybersecurity infrastructure, tasked with identifying malicious activities across increasingly diverse and complex network environments. Contemporary network ecosystems encompass heterogeneous domains including cloud platforms, industrial IoT systems, containerized infrastructure, enterprise networks, consumer IoT devices, and hybrid architectures. Each domain exhibits distinct traffic characteristics, attack patterns, and operational constraints, posing significant challenges for unified detection frameworks.

Recent advances in deep learning have demonstrated promising results for intrusion detection, with techniques ranging from recurrent neural networks~\cite{reference1} to graph neural networks~\cite{reference2} and continuous-depth models~\cite{reference3}. However, existing approaches suffer from three fundamental limitations that hinder their deployment in safety-critical security contexts:

\textbf{(1) Lack of Cross-Domain Generalization:} Most deep learning NIDS models are designed and optimized for specific network domains, requiring extensive retraining or architectural modifications when applied to new environments. This domain specificity limits practical applicability in heterogeneous infrastructure.

\textbf{(2) Inadequate Uncertainty Quantification:} Production security systems require not only accurate predictions but also calibrated confidence estimates to support risk-aware decision making. Standard neural networks produce overconfident predictions, while existing uncertainty quantification methods like Monte Carlo dropout or deep ensembles incur substantial computational overhead during inference.

\textbf{(3) Vulnerability to Adversarial Attacks:} Adversarial examples crafted to evade detection pose serious threats to ML-based NIDS. While empirical defenses exist, few approaches provide certified robustness guarantees that bound the worst-case attack success rate.

To address these limitations, we propose SDE-TGNN, a novel architecture that models network traffic evolution as a stochastic differential equation over temporal graph structures. Our key insight is that network states evolve continuously in time according to both deterministic dynamics (representing normal operational patterns and systematic attacks) and stochastic fluctuations (capturing traffic variability and unpredictable anomalies). By explicitly modeling these components via learned drift and diffusion functions, SDE-TGNN achieves three critical advantages:

\textbf{First}, the continuous-time formulation provides a natural inductive bias for temporal patterns that generalizes across domains with different sampling rates and temporal scales. The multi-scale temporal fusion mechanism enables simultaneous modeling of short-term anomalies (e.g., port scans, DDoS bursts) and long-term attack campaigns (e.g., APT reconnaissance, data exfiltration).

\textbf{Second}, uncertainty quantification emerges naturally from the stochastic formulation. Rather than expensive Monte Carlo sampling, we propagate the first two moments (mean and covariance) of the state distribution analytically using the Fokker-Planck equation. This moment propagation approach achieves $O(d^2)$ complexity in state dimension $d$, orders of magnitude faster than sampling-based methods while providing principled uncertainty estimates.

\textbf{Third}, the inherent stochasticity in the SDE framework enables certified adversarial robustness via randomized smoothing. By introducing controlled Gaussian noise and leveraging concentration inequalities, we derive provable lower bounds on attack perturbation budgets required to change predictions, providing defense guarantees against $\ell_2$-bounded adversaries.

\subsection{Contributions}

The main contributions of this work are:

\begin{enumerate}[leftmargin=*]
\item \textbf{SDE-Based Temporal Dynamics:} We introduce a novel architecture that models continuous-time network state evolution through learned drift $f_\theta(h_t, \G_t, t)$ and diffusion $\sigma_\phi(h_t, t)$ functions, capturing both deterministic trends and stochastic variability in traffic patterns.

\item \textbf{Fokker-Planck Uncertainty Propagation:} Instead of Monte Carlo sampling, we derive analytical moment propagation equations from the Fokker-Planck formulation, enabling efficient $O(d^2)$ uncertainty quantification with calibrated confidence estimates.

\item \textbf{Multi-Scale Temporal Fusion:} We integrate SDE trajectories at multiple temporal scales (short, medium, long-term) through a learned fusion mechanism, enabling simultaneous detection of diverse attack types with different temporal signatures.

\item \textbf{Certified Adversarial Robustness:} We leverage the stochastic nature of the model to provide certified accuracy guarantees via randomized smoothing, deriving provable lower bounds on adversarial perturbation budgets.

\item \textbf{Comprehensive Multi-Domain Evaluation:} We conduct extensive experiments on six heterogeneous network security datasets (Microsoft Azure Cloud, Edge-IIoTset, Kubernetes vs Docker, CSE-CIC-IDS2018, UNB-CIC-IoT2023, UNSW-NB15) totaling 53+ million samples, demonstrating state-of-the-art detection performance, calibrated uncertainty, and certified robustness without domain-specific modifications.

\item \textbf{Open-Source Implementation:} We release a comprehensive, modular implementation including data preprocessing pipelines, model architectures, training procedures, and evaluation protocols to facilitate reproducibility and future research.
\end{enumerate}

The remainder of this paper is organized as follows: Section~II reviews related work; Section~III formulates the problem; Section~IV presents the SDE-TGNN architecture; Section~V describes the Fokker-Planck uncertainty propagation; Section~VI details the multi-scale temporal fusion; Section~VII presents certified robustness guarantees; Section~VIII describes experimental setup and datasets; Section~IX presents results and analysis; Section~X discusses implications and limitations; Section~XI concludes.

\section{Related Work}

\subsection{Deep Learning for Network Intrusion Detection}

Early deep learning approaches to NIDS focused on recurrent architectures, with LSTM and BiLSTM networks demonstrating strong performance on sequential packet data~\cite{ref_lstm}. Convolutional neural networks have been applied to extract spatial features from network traffic representations~\cite{ref_cnn}. Hybrid CNN-LSTM architectures combine spatial feature extraction with temporal modeling~\cite{ref_cnnlstm}. However, these approaches treat time as discrete steps and struggle with long-range temporal dependencies.

More recent work has explored transformer architectures~\cite{ref_transformer} for modeling attention patterns across traffic sequences. While transformers handle long sequences effectively, their quadratic computational complexity limits scalability. Attention mechanisms have also been integrated with graph neural networks to model network topology~\cite{ref_gat}, but these approaches typically operate on static or discrete-time graph snapshots.

\subsection{Graph Neural Networks for Network Security}

Graph neural networks provide a natural framework for representing network topology and communication patterns. GraphSAGE~\cite{ref_graphsage} and GAT~\cite{ref_gat} have been applied to network traffic graphs, learning node embeddings that capture structural properties. Temporal graph networks~\cite{ref_tgn} extend this paradigm to evolving graphs through recurrent aggregation of temporal neighborhoods.

However, existing temporal GNN approaches use discrete-time formulations that are sensitive to sampling rates and struggle to model continuous temporal dynamics. They also lack principled uncertainty quantification, producing point predictions without confidence estimates. Our SDE-TGNN framework addresses both limitations through continuous-time stochastic modeling and analytical moment propagation.

\subsection{Continuous-Depth Models}

Neural Ordinary Differential Equations (Neural ODEs)~\cite{ref_node} parameterize hidden states as solutions to ODEs, enabling continuous-depth architectures with memory-efficient training via the adjoint method. Neural ODEs have been applied to time series modeling~\cite{ref_latentode} and dynamics learning~\cite{ref_augmentedode}. Graph Neural ODEs~\cite{ref_gnode} extend this framework to graphs, modeling node feature evolution through continuous-time GNN dynamics.

While Neural ODEs provide continuous-time modeling, they are inherently deterministic and cannot quantify uncertainty arising from stochastic traffic variability. Stochastic differential equations address this limitation by introducing diffusion terms that model random fluctuations. Prior work on Neural SDEs~\cite{ref_neuralsde} has focused on time series forecasting and generative modeling, but has not been applied to graph-structured security data with principled uncertainty quantification.

\subsection{Uncertainty Quantification in Deep Learning}

Bayesian neural networks~\cite{ref_bnn} provide a principled framework for uncertainty quantification through posterior distributions over weights. However, exact Bayesian inference is intractable, requiring approximations like variational inference~\cite{ref_vi} or Monte Carlo dropout~\cite{ref_mcdropout}. Deep ensembles~\cite{ref_ensemble} offer an alternative approach, training multiple models with different initializations. While effective, these methods incur substantial computational overhead during inference (proportional to the number of samples or ensemble members).

Our Fokker-Planck-based approach propagates uncertainty analytically through the forward dynamics, avoiding repeated sampling. This is related to moment propagation methods in Kalman filtering~\cite{ref_kalman} and unscented transforms~\cite{ref_unscented}, but extends these techniques to high-dimensional neural network states with learned stochastic dynamics.

\subsection{Adversarial Robustness}

Adversarial training~\cite{ref_advtrain} improves empirical robustness by augmenting training data with adversarial examples. However, empirical defenses often fail against stronger attacks. Certified defense methods provide provable guarantees, including convex relaxation~\cite{ref_convex}, linear bound propagation~\cite{ref_crown}, and randomized smoothing~\cite{ref_randomized}.

Randomized smoothing constructs a smoothed classifier by averaging predictions over Gaussian noise, deriving certified radii via concentration inequalities. This approach has been applied to image classification~\cite{ref_cohen} and time series~\cite{ref_tsrobust}, but not to graph-structured network traffic data with continuous-time dynamics. Our work demonstrates that the inherent stochasticity in the SDE formulation naturally enables randomized smoothing-based certified robustness.

\section{Problem Formulation}

\subsection{Network Traffic as Temporal Graphs}

Consider a network at time $t$ represented as a temporal graph $\G_t = (\V_t, \mathcal{E}_t, \X_t)$, where:
\begin{itemize}[leftmargin=*]
\item $\V_t$ is the set of network entities (hosts, services, devices)
\item $\mathcal{E}_t \subseteq \V_t \times \V_t$ represents communication edges
\item $\X_t \in \R^{|\V_t| \times d_{\text{in}}}$ contains node features (traffic statistics, protocol information, payload characteristics)
\end{itemize}

Network traffic evolves continuously, with $\G_t$ changing as connections form/terminate and features update. We observe a sequence of graph snapshots $\{\G_{t_1}, \G_{t_2}, \ldots, \G_{t_T}\}$ sampled at discrete times.

\subsection{Intrusion Detection Task}

For each observation at time $t_i$, the goal is to predict:
\begin{equation}
y_i \in \{0, 1, \ldots, C-1\}
\end{equation}
where $y_i = 0$ denotes benign traffic and $y_i > 0$ indicates attack class $y_i$.

Beyond point predictions, production security systems require:
\begin{enumerate}[leftmargin=*]
\item \textbf{Uncertainty estimates:} Confidence scores $p(y_i | \G_{1:t_i})$ that quantify prediction reliability
\item \textbf{Robustness guarantees:} Certified bounds on adversarial perturbation budgets
\item \textbf{Cross-domain generalization:} Performance across diverse network domains without retraining
\end{enumerate}

\subsection{Multi-Domain Setting}

We consider $K$ heterogeneous network domains $\{\D_1, \ldots, \D_K\}$, each with domain-specific:
\begin{itemize}[leftmargin=*]
\item Feature distributions: Traffic statistics vary across cloud, IoT, enterprise, etc.
\item Attack patterns: Different domains face domain-specific threats
\item Temporal characteristics: Sampling rates and traffic dynamics differ
\end{itemize}

The challenge is to learn a unified model $f_\theta$ that achieves strong performance across all domains:
\begin{equation}
\theta^* = \argmin_\theta \sum_{k=1}^K \E_{(\G, y) \sim \D_k} [\mathcal{L}(f_\theta(\G), y)]
\end{equation}
without domain-specific adaptations or fine-tuning.

\section{SDE-TGNN Architecture}

\subsection{Overview}

The SDE-TGNN architecture consists of four main components:

\begin{enumerate}[leftmargin=*]
\item \textbf{Feature Embedding:} Maps heterogeneous input features to a unified representation space
\item \textbf{Graph Attention Layers:} Learns structural representations via multi-head attention over temporal neighborhoods
\item \textbf{SDE Temporal Dynamics:} Models continuous-time state evolution through learned drift and diffusion
\item \textbf{Multi-Scale Fusion \& Classification:} Integrates trajectories at multiple temporal scales for final prediction
\end{enumerate}

\subsection{Feature Embedding}

Given input node features $\X_t \in \R^{n \times d_{\text{in}}}$, we apply a learned embedding:
\begin{equation}
\mathbf{H}^{(0)}_t = \text{LayerNorm}(\X_t \mathbf{W}_{\text{emb}} + \mathbf{b}_{\text{emb}})
\end{equation}
where $\mathbf{W}_{\text{emb}} \in \R^{d_{\text{in}} \times D}$ projects to hidden dimension $D$. Layer normalization ensures stable training across domains with different feature scales.

\subsection{Temporal Graph Attention}

We apply $L$ graph attention layers~\cite{ref_gat} to learn node representations that incorporate structural information:
\begin{equation}
\mathbf{H}^{(\ell+1)}_t = \text{GAT}(\mathbf{H}^{(\ell)}_t, \mathcal{E}_t) = \bigoplus_{k=1}^{K_h} \sigma\left(\sum_{j \in \mathcal{N}_i} \alpha_{ij}^k \mathbf{W}_k^{(\ell)} \mathbf{h}_j^{(\ell)}\right)
\end{equation}
where $K_h$ is the number of attention heads, $\mathcal{N}_i$ denotes neighbors of node $i$, and attention coefficients $\alpha_{ij}^k$ are computed via:
\begin{equation}
\alpha_{ij}^k = \frac{\exp(\text{LeakyReLU}(\mathbf{a}_k^T [\mathbf{W}_k \mathbf{h}_i \| \mathbf{W}_k \mathbf{h}_j]))}{\sum_{j' \in \mathcal{N}_i} \exp(\text{LeakyReLU}(\mathbf{a}_k^T [\mathbf{W}_k \mathbf{h}_i \| \mathbf{W}_k \mathbf{h}_{j'}]))}
\end{equation}

After $L$ layers, we obtain graph-aware node representations $\mathbf{H}^{(L)}_t \in \R^{n \times D}$.

\subsection{State Projection}

To enable efficient SDE integration, we project the high-dimensional graph representations to a lower-dimensional state space:
\begin{equation}
\mathbf{h}_t = \text{MeanPool}(\mathbf{H}^{(L)}_t) \mathbf{W}_{\text{proj}} \in \R^d
\end{equation}
where $d \ll D$ is the SDE state dimension. Mean pooling aggregates node-level features into a graph-level representation, and $\mathbf{W}_{\text{proj}} \in \R^{D \times d}$ performs dimensionality reduction.

\subsection{Stochastic Differential Equation Dynamics}

We model the continuous-time evolution of the network state $\mathbf{h}_t$ via the SDE:
\begin{equation}
d\mathbf{h}_t = f_\theta(\mathbf{h}_t, \G_t, t) dt + \sigma_\phi(\mathbf{h}_t, t) d\mathbf{W}_t
\label{eq:sde}
\end{equation}
where:
\begin{itemize}[leftmargin=*]
\item $f_\theta: \R^d \times \G \times \R^+ \to \R^d$ is the \textbf{drift function} capturing deterministic dynamics
\item $\sigma_\phi: \R^d \times \R^+ \to \R^{d \times d}$ is the \textbf{diffusion function} modeling stochastic variability
\item $\mathbf{W}_t$ is a standard $d$-dimensional Brownian motion
\end{itemize}

\subsubsection{Drift Network}

The drift network parameterizes deterministic state evolution:
\begin{equation}
f_\theta(\mathbf{h}_t, \G_t, t) = \text{MLP}_{\text{drift}}([\mathbf{h}_t \| \phi(\G_t) \| \psi(t)])
\end{equation}
where:
\begin{itemize}[leftmargin=*]
\item $\phi(\G_t)$ extracts graph-level statistics (edge density, clustering coefficient, etc.)
\item $\psi(t) = [\sin(\omega_1 t), \cos(\omega_1 t), \ldots, \sin(\omega_m t), \cos(\omega_m t)]$ provides time encoding
\item $\text{MLP}_{\text{drift}}$ is a multi-layer perceptron with $[\cdot \| \cdot \| \cdot]$ denoting concatenation
\end{itemize}

\subsubsection{Diffusion Network}

The diffusion network models state-dependent and time-varying noise:
\begin{equation}
\sigma_\phi(\mathbf{h}_t, t) = \text{Softplus}(\text{MLP}_{\text{diff}}([\mathbf{h}_t \| \psi(t)])) \cdot \mathbf{I}_d
\end{equation}
We use a diagonal diffusion matrix (scaled identity) for computational efficiency. The softplus activation $\text{Softplus}(x) = \log(1 + e^x)$ ensures positive-definite diffusion coefficients.

\subsection{Multi-Scale SDE Integration}

To capture temporal patterns at different scales, we integrate the SDE over three time horizons:
\begin{align}
\mathbf{h}_t^{(s)} &= \text{SDE-Solve}(\mathbf{h}_0, [0, \tau_s], f_\theta, \sigma_\phi) \quad \text{(short-term)} \\
\mathbf{h}_t^{(m)} &= \text{SDE-Solve}(\mathbf{h}_0, [0, \tau_m], f_\theta, \sigma_\phi) \quad \text{(medium-term)} \\
\mathbf{h}_t^{(l)} &= \text{SDE-Solve}(\mathbf{h}_0, [0, \tau_l], f_\theta, \sigma_\phi) \quad \text{(long-term)}
\end{align}
where $\tau_s < \tau_m < \tau_l$ are time horizons (e.g., $\tau_s = 0.5$s, $\tau_m = 2$s, $\tau_l = 10$s).

For numerical integration, we employ the Euler-Maruyama scheme:
\begin{equation}
\mathbf{h}_{t+\Delta t} = \mathbf{h}_t + f_\theta(\mathbf{h}_t, \G_t, t) \Delta t + \sigma_\phi(\mathbf{h}_t, t) \sqrt{\Delta t} \boldsymbol{\epsilon}_t
\end{equation}
where $\boldsymbol{\epsilon}_t \sim \N(0, \mathbf{I}_d)$ and $\Delta t$ is the integration time step.

\subsection{Multi-Scale Fusion}

We fuse the multi-scale trajectories via a learned attention mechanism:
\begin{equation}
\mathbf{h}_{\text{fused}} = \sum_{k \in \{s,m,l\}} \beta_k \mathbf{h}_t^{(k)}
\end{equation}
where attention weights $\boldsymbol{\beta} = \text{Softmax}(\mathbf{W}_{\text{attn}} [\mathbf{h}_t^{(s)} \| \mathbf{h}_t^{(m)} \| \mathbf{h}_t^{(l)}])$ are learned from the concatenated multi-scale states.

\subsection{Classification Head}

The fused representation is passed through a classification MLP:
\begin{equation}
\mathbf{z} = \text{MLP}_{\text{cls}}(\mathbf{h}_{\text{fused}}) \in \R^C
\end{equation}
producing logits $\mathbf{z}$ for $C$ classes. We apply softmax to obtain predictive probabilities:
\begin{equation}
p(y = c | \G) = \frac{\exp(z_c)}{\sum_{c'=1}^C \exp(z_{c'})}
\end{equation}

\section{Fokker-Planck Uncertainty Propagation}

A key advantage of the SDE formulation is principled uncertainty quantification. Rather than expensive Monte Carlo sampling, we propagate the distribution $p(\mathbf{h}_t | \mathbf{h}_0)$ analytically via the Fokker-Planck equation.

\subsection{Fokker-Planck Equation}

The probability density $p(\mathbf{h}, t)$ of the SDE state evolves according to:
\begin{equation}
\frac{\partial p}{\partial t} = -\nabla \cdot (f_\theta p) + \frac{1}{2} \Tr(\sigma_\phi \sigma_\phi^T \nabla \nabla^T p)
\label{eq:fp}
\end{equation}
This partial differential equation is intractable for high-dimensional neural network states. Instead, we propagate the first two moments.

\subsection{Moment Propagation}

Let $\boldsymbol{\mu}_t = \E[\mathbf{h}_t]$ and $\boldsymbol{\Sigma}_t = \text{Cov}[\mathbf{h}_t]$. Taking expectations of the SDE~\eqref{eq:sde}, we obtain:
\begin{align}
\frac{d\boldsymbol{\mu}_t}{dt} &= \E[f_\theta(\mathbf{h}_t, \G_t, t)] \label{eq:moment1} \\
\frac{d\boldsymbol{\Sigma}_t}{dt} &= \E[J_{f_\theta}(\mathbf{h}_t)] \boldsymbol{\Sigma}_t + \boldsymbol{\Sigma}_t \E[J_{f_\theta}(\mathbf{h}_t)]^T + \E[\sigma_\phi(\mathbf{h}_t, t) \sigma_\phi(\mathbf{h}_t, t)^T] \label{eq:moment2}
\end{align}
where $J_{f_\theta}$ denotes the Jacobian of the drift function.

\subsection{Linearization Approximation}

To make moment propagation tractable, we linearize the drift around the current mean:
\begin{equation}
f_\theta(\mathbf{h}_t, \G_t, t) \approx f_\theta(\boldsymbol{\mu}_t, \G_t, t) + J_{f_\theta}(\boldsymbol{\mu}_t) (\mathbf{h}_t - \boldsymbol{\mu}_t)
\end{equation}
Substituting into~\eqref{eq:moment1}--\eqref{eq:moment2} and simplifying:
\begin{align}
\frac{d\boldsymbol{\mu}_t}{dt} &= f_\theta(\boldsymbol{\mu}_t, \G_t, t) \label{eq:mean_ode} \\
\frac{d\boldsymbol{\Sigma}_t}{dt} &= J_{f_\theta}(\boldsymbol{\mu}_t) \boldsymbol{\Sigma}_t + \boldsymbol{\Sigma}_t J_{f_\theta}(\boldsymbol{\mu}_t)^T + Q_t \label{eq:cov_ode}
\end{align}
where $Q_t = \sigma_\phi(\boldsymbol{\mu}_t, t) \sigma_\phi(\boldsymbol{\mu}_t, t)^T$ is the process noise covariance.

\subsection{Efficient Implementation}

Equations~\eqref{eq:mean_ode}--\eqref{eq:cov_ode} form a coupled ODE system in $O(d + d^2)$ variables. We integrate this system jointly with the SDE trajectories:
\begin{itemize}[leftmargin=*]
\item Forward pass: Solve~\eqref{eq:mean_ode}--\eqref{eq:cov_ode} from $t=0$ to $t=T$ to obtain $\boldsymbol{\mu}_T, \boldsymbol{\Sigma}_T$
\item Uncertainty quantification: Predictive variance for dimension $i$ is $[\boldsymbol{\Sigma}_T]_{ii}$
\item Total uncertainty: $\text{trace}(\boldsymbol{\Sigma}_T) = \sum_{i=1}^d [\boldsymbol{\Sigma}_T]_{ii}$
\end{itemize}

Jacobian $J_{f_\theta}$ is computed efficiently via automatic differentiation. The overall complexity is $O(d^2)$ per integration step, far more efficient than Monte Carlo sampling which requires $O(Nd)$ for $N$ samples with typical $N \gg d$.

\subsection{Uncertainty-Aware Loss}

We incorporate uncertainty into the training objective via a Bayesian risk-sensitive loss:
\begin{equation}
\mathcal{L}_{\text{unc}} = \E_{p(\mathbf{h}_T | \mathbf{h}_0)} [\mathcal{L}_{\text{CE}}(\mathbf{z}(\mathbf{h}_T), y)]
\end{equation}
where $\mathcal{L}_{\text{CE}}$ is cross-entropy loss. Under the Gaussian approximation $p(\mathbf{h}_T | \mathbf{h}_0) \approx \N(\boldsymbol{\mu}_T, \boldsymbol{\Sigma}_T)$, we use a second-order Taylor expansion:
\begin{equation}
\mathcal{L}_{\text{unc}} \approx \mathcal{L}_{\text{CE}}(\mathbf{z}(\boldsymbol{\mu}_T), y) + \frac{1}{2} \Tr(\mathbf{H}_{\mathbf{z}} \boldsymbol{\Sigma}_T)
\end{equation}
where $\mathbf{H}_{\mathbf{z}}$ is the Hessian of the classification head. The second term penalizes high uncertainty, encouraging the model to reduce prediction variance.

\section{Multi-Scale Temporal Fusion}

Network attacks exhibit diverse temporal signatures:
\begin{itemize}[leftmargin=*]
\item \textbf{Short-term:} Port scans, DoS bursts (seconds)
\item \textbf{Medium-term:} Brute force attempts, reconnaissance (minutes)
\item \textbf{Long-term:} APT campaigns, data exfiltration (hours/days)
\end{itemize}

Single-scale models struggle to capture this spectrum. Our multi-scale fusion mechanism addresses this by integrating SDE trajectories at three temporal horizons.

\subsection{Scale-Specific Dynamics}

At scale $k \in \{s, m, l\}$ with time horizon $\tau_k$, we solve:
\begin{equation}
\mathbf{h}_{\tau_k}^{(k)} = \mathbf{h}_0 + \int_0^{\tau_k} f_\theta(\mathbf{h}_t, \G_t, t) dt + \int_0^{\tau_k} \sigma_\phi(\mathbf{h}_t, t) d\mathbf{W}_t
\end{equation}

Different scales emphasize different frequency components of the dynamics. Short-term trajectories ($\tau_s \approx 0.5$s) capture rapid fluctuations, while long-term trajectories ($\tau_l \approx 10$s) integrate over slower trends.

\subsection{Attention-Based Fusion}

We compute scale-importance weights via self-attention:
\begin{align}
\mathbf{Q} &= \mathbf{W}_Q [\mathbf{h}_{\tau_s}^{(s)} \| \mathbf{h}_{\tau_m}^{(m)} \| \mathbf{h}_{\tau_l}^{(l)}] \\
\mathbf{K} &= \mathbf{W}_K [\mathbf{h}_{\tau_s}^{(s)} \| \mathbf{h}_{\tau_m}^{(m)} \| \mathbf{h}_{\tau_l}^{(l)}] \\
\boldsymbol{\beta} &= \text{Softmax}(\mathbf{Q}^T \mathbf{K} / \sqrt{3d})
\end{align}
and fuse:
\begin{equation}
\mathbf{h}_{\text{fused}} = \sum_{k \in \{s,m,l\}} \beta_k \mathbf{h}_{\tau_k}^{(k)}
\end{equation}

This allows the model to adaptively weight scales based on input characteristics. For instance, DDoS attacks may receive higher weight on short-term scales, while APT reconnaissance emphasizes long-term scales.

\subsection{Scale-Specific Uncertainty}

Each scale maintains separate moment statistics $(\boldsymbol{\mu}_{\tau_k}^{(k)}, \boldsymbol{\Sigma}_{\tau_k}^{(k)})$. The fused uncertainty is:
\begin{equation}
\boldsymbol{\Sigma}_{\text{fused}} = \sum_{k} \beta_k^2 \boldsymbol{\Sigma}_{\tau_k}^{(k)} + \sum_k \beta_k (\boldsymbol{\mu}_{\tau_k}^{(k)} - \boldsymbol{\mu}_{\text{fused}})(\boldsymbol{\mu}_{\tau_k}^{(k)} - \boldsymbol{\mu}_{\text{fused}})^T
\end{equation}
capturing both within-scale and between-scale variability.

\section{Certified Adversarial Robustness}

The stochastic nature of SDE-TGNN enables certified robustness guarantees via randomized smoothing~\cite{ref_cohen}.

\subsection{Randomized Smoothing Framework}

Given a base classifier $f(\mathbf{x})$ and noise distribution $\eta \sim \N(0, \sigma^2 \mathbf{I})$, the smoothed classifier is:
\begin{equation}
g(\mathbf{x}) = \argmax_c \Pr[f(\mathbf{x} + \eta) = c]
\end{equation}

\begin{theorem}[Cohen et al.~\cite{ref_cohen}]
If $g(\mathbf{x}) = c_A$ with $\Pr[f(\mathbf{x} + \eta) = c_A] \geq p_A$ and $\max_{c \neq c_A} \Pr[f(\mathbf{x} + \eta) = c] \leq p_B$ where $p_A > p_B$, then $g$ is constant within $\ell_2$ ball of radius:
\begin{equation}
R = \frac{\sigma}{2} (\Phi^{-1}(p_A) - \Phi^{-1}(p_B))
\end{equation}
\end{theorem}

\subsection{Integration with SDE-TGNN}

Our SDE formulation provides inherent stochasticity through the diffusion term. We leverage this in two ways:

\textbf{(1) Training-time smoothing:} The diffusion $\sigma_\phi$ introduces controlled noise during forward propagation, acting as implicit data augmentation. This improves robustness without explicit adversarial training.

\textbf{(2) Inference-time certification:} For critical predictions, we sample $N$ trajectories from the SDE, compute class probabilities $\{p_1, \ldots, p_N\}$, estimate $p_A$ and $p_B$ via Monte Carlo, and apply the certified radius formula.

\subsection{Implementation}

For certification, we:
\begin{enumerate}[leftmargin=*]
\item Sample $N=10000$ SDE trajectories with initial state $\mathbf{h}_0 \sim \N(\bar{\mathbf{h}}_0, \sigma_{\text{cert}}^2 \mathbf{I})$
\item Compute predictions $\{c_1, \ldots, c_N\}$
\item Estimate class probabilities: $\hat{p}_c = \frac{1}{N} \sum_{i=1}^N \mathbb{I}[c_i = c]$
\item Compute lower confidence bound on $p_A$ and upper bound on $p_B$ using Clopper-Pearson intervals
\item Calculate certified radius $R$
\end{enumerate}

The certification noise level $\sigma_{\text{cert}}$ is a hyperparameter balancing accuracy and certified radius. Larger $\sigma_{\text{cert}}$ provides larger certified radii but may reduce clean accuracy.

\section{Experimental Setup}

\subsection{Datasets}

We evaluate on six heterogeneous network security datasets:

\begin{table}[t]
\centering
\caption{Multi-Domain Network Intrusion Detection Datasets}
\label{tab:datasets}
\small
\begin{tabular}{@{}lcccc@{}}
\toprule
\textbf{Dataset} & \textbf{Domain} & \textbf{Samples} & \textbf{Classes} & \textbf{Features} \\
\midrule
Microsoft Azure & Cloud & 1.0M & 2 & 81 \\
Edge-IIoTset & Industrial IoT & 157K & 15 & 61 \\
K8s vs Docker & Container & 50K & 5 & 47 \\
CSE-CIC-IDS2018 & Enterprise & 16.2M & 15 & 79 \\
UNB-CIC-IoT2023 & IoT Network & 33.7M & 34 & 46 \\
UNSW-NB15 & Hybrid & 2.5M & 10 & 49 \\
\midrule
\textbf{Total} & -- & \textbf{53.6M} & -- & -- \\
\bottomrule
\end{tabular}
\end{table}

\textbf{Microsoft Azure Cloud:} Cloud-based malware prediction with system configuration and security features.

\textbf{Edge-IIoTset:} Industrial IoT sensors and actuators with 14 attack types (DDoS, injection, backdoor, etc.).

\textbf{Kubernetes vs Docker:} Container orchestration security with pod-level traffic statistics.

\textbf{CSE-CIC-IDS2018:} Enterprise network with diverse attacks (brute force, DoS, web attacks, infiltration, botnet).

\textbf{UNB-CIC-IoT2023:} Comprehensive IoT dataset with 33 attack classes across smart home and industrial devices.

\textbf{UNSW-NB15:} Hybrid network combining normal activities and contemporary attack vectors.

\subsection{Data Preprocessing}

For each dataset, we:
\begin{enumerate}[leftmargin=*]
\item \textbf{Feature harmonization:} Map raw features to standardized categories (packet statistics, flow duration, protocol flags, etc.)
\item \textbf{Normalization:} Apply z-score normalization per dataset to handle different scales
\item \textbf{Graph construction:} Build $k$-NN graphs ($k=10$) based on feature similarity to capture flow relationships
\item \textbf{Train/val/test split:} 70\%/15\%/15\% stratified split
\end{enumerate}

\subsection{Baseline Models}

We compare against eight baselines:

\textbf{Classical ML:}
\begin{itemize}[leftmargin=*]
\item Random Forest (RF): 200 trees, max depth 20
\item XGBoost: 200 estimators, learning rate 0.1
\end{itemize}

\textbf{Deep Learning:}
\begin{itemize}[leftmargin=*]
\item BiLSTM: 2 layers, 256 hidden units
\item CNN-BiLSTM: 3 conv layers + 2 BiLSTM layers
\end{itemize}

\textbf{Graph Neural Networks:}
\begin{itemize}[leftmargin=*]
\item GraphSAGE: 3 layers, mean aggregation, 256 hidden
\end{itemize}

\textbf{Continuous Models:}
\begin{itemize}[leftmargin=*]
\item Neural ODE: 4 ODE function layers, adjoint method, 256 hidden
\end{itemize}

\textbf{Uncertainty Quantification:}
\begin{itemize}[leftmargin=*]
\item MC Dropout: BiLSTM with dropout 0.3, 50 samples at inference
\item Deep Ensemble: 5 BiLSTM models with different seeds
\end{itemize}

\subsection{Training Configuration}

\textbf{Model Hyperparameters:}
\begin{itemize}[leftmargin=*]
\item Graph attention: $L=4$ layers, $D=256$ hidden, $K_h=8$ heads
\item SDE state: $d=64$ dimensions
\item Temporal scales: $\tau_s=0.5$s, $\tau_m=2$s, $\tau_l=10$s
\item SDE solver: Euler-Maruyama, $\Delta t = 0.01$
\end{itemize}

\textbf{Training:}
\begin{itemize}[leftmargin=*]
\item Optimizer: AdamW, lr=$2 \times 10^{-4}$, weight decay $10^{-4}$
\item Batch size: 128
\item Max epochs: 100
\item Early stopping: patience 15 epochs on validation loss
\item Loss: $\mathcal{L} = \mathcal{L}_{\text{CE}} + 0.1 \mathcal{L}_{\text{unc}}$
\end{itemize}

\subsection{Evaluation Metrics}

\textbf{Detection Performance:}
\begin{itemize}[leftmargin=*]
\item Accuracy, Precision, Recall, F1 score (macro and weighted)
\item AUC-ROC (one-vs-rest for multi-class)
\end{itemize}

\textbf{Uncertainty Calibration:}
\begin{itemize}[leftmargin=*]
\item Expected Calibration Error (ECE) with 15 bins
\item Maximum Calibration Error (MCE)
\item Brier score
\item Reliability diagrams
\end{itemize}

\textbf{Adversarial Robustness:}
\begin{itemize}[leftmargin=*]
\item FGSM accuracy at $\epsilon \in \{0.01, 0.05, 0.1\}$
\item PGD accuracy (20 iterations, step size 0.01)
\item Certified accuracy at radii $R \in \{0.1, 0.25, 0.5\}$
\end{itemize}

All experiments use 5-fold cross-validation with different random seeds. We report mean and standard deviation across folds.

\section{Results and Analysis}

\subsection{Detection Performance}

\begin{table*}[t]
\centering
\caption{Detection Performance Across Six Datasets (Weighted F1 Score \%)}
\label{tab:detection_performance}
\small
\begin{tabular}{@{}lccccccc@{}}
\toprule
\textbf{Model} & \textbf{Azure} & \textbf{IIoT} & \textbf{K8s} & \textbf{CIC-IDS18} & \textbf{IoT-2023} & \textbf{UNSW} & \textbf{Avg} \\
\midrule
Random Forest & 87.2 $\pm$ 0.5 & 92.1 $\pm$ 0.8 & 84.3 $\pm$ 1.2 & 89.5 $\pm$ 0.6 & 85.7 $\pm$ 0.9 & 88.1 $\pm$ 0.7 & 87.8 \\
XGBoost & 89.5 $\pm$ 0.6 & 93.7 $\pm$ 0.7 & 86.8 $\pm$ 1.0 & 91.2 $\pm$ 0.5 & 87.4 $\pm$ 0.8 & 89.9 $\pm$ 0.6 & 89.8 \\
BiLSTM & 92.1 $\pm$ 0.4 & 94.5 $\pm$ 0.6 & 88.2 $\pm$ 0.9 & 93.8 $\pm$ 0.4 & 90.2 $\pm$ 0.7 & 91.7 $\pm$ 0.5 & 91.8 \\
CNN-BiLSTM & 93.8 $\pm$ 0.5 & 95.3 $\pm$ 0.5 & 89.7 $\pm$ 0.8 & 94.6 $\pm$ 0.4 & 91.5 $\pm$ 0.6 & 92.8 $\pm$ 0.5 & 93.0 \\
GraphSAGE & 94.2 $\pm$ 0.4 & 96.1 $\pm$ 0.4 & 90.5 $\pm$ 0.7 & 95.3 $\pm$ 0.3 & 92.8 $\pm$ 0.5 & 93.6 $\pm$ 0.4 & 93.8 \\
Neural ODE & 95.1 $\pm$ 0.4 & 96.8 $\pm$ 0.3 & 91.2 $\pm$ 0.6 & 96.0 $\pm$ 0.3 & 93.5 $\pm$ 0.5 & 94.3 $\pm$ 0.4 & 94.5 \\
MC Dropout & 92.5 $\pm$ 0.6 & 94.8 $\pm$ 0.7 & 88.5 $\pm$ 1.1 & 94.1 $\pm$ 0.5 & 90.6 $\pm$ 0.8 & 92.0 $\pm$ 0.6 & 92.1 \\
Deep Ensemble & 94.8 $\pm$ 0.3 & 96.5 $\pm$ 0.4 & 90.8 $\pm$ 0.5 & 95.7 $\pm$ 0.3 & 93.1 $\pm$ 0.4 & 94.0 $\pm$ 0.4 & 94.2 \\
\midrule
\textbf{SDE-TGNN} & \textbf{98.3 $\pm$ 0.2} & \textbf{99.1 $\pm$ 0.2} & \textbf{97.5 $\pm$ 0.3} & \textbf{99.4 $\pm$ 0.1} & \textbf{98.9 $\pm$ 0.2} & \textbf{98.8 $\pm$ 0.2} & \textbf{98.7} \\
\bottomrule
\end{tabular}
\end{table*}

Table~\ref{tab:detection_performance} shows SDE-TGNN achieves state-of-the-art performance across all six datasets, with an average weighted F1 score of 98.7\%, outperforming the next-best baseline (Neural ODE at 94.5\%) by 4.2 percentage points. Performance gains are consistent across domains:

\begin{itemize}[leftmargin=*]
\item \textbf{Cloud (Azure):} 98.3\% F1, +3.2\% over Neural ODE
\item \textbf{Industrial IoT:} 99.1\% F1, +2.3\% over Neural ODE
\item \textbf{Container:} 97.5\% F1, +6.3\% over Neural ODE
\item \textbf{Enterprise:} 99.4\% F1, +3.4\% over Neural ODE
\item \textbf{IoT:} 98.9\% F1, +5.4\% over Neural ODE
\item \textbf{Hybrid:} 98.8\% F1, +4.5\% over Neural ODE
\end{itemize}

The largest improvements occur on container (K8s) and IoT datasets, which exhibit high temporal variability that the stochastic modeling captures effectively.

\subsection{Uncertainty Calibration}

\begin{table}[t]
\centering
\caption{Uncertainty Calibration Metrics (Lower is Better)}
\label{tab:calibration}
\small
\begin{tabular}{@{}lccc@{}}
\toprule
\textbf{Model} & \textbf{ECE} $\downarrow$ & \textbf{MCE} $\downarrow$ & \textbf{Brier} $\downarrow$ \\
\midrule
BiLSTM & 0.187 & 0.324 & 0.142 \\
MC Dropout & 0.092 & 0.176 & 0.089 \\
Deep Ensemble & 0.078 & 0.145 & 0.071 \\
\textbf{SDE-TGNN} & \textbf{0.042} & \textbf{0.098} & \textbf{0.038} \\
\bottomrule
\end{tabular}
\end{table}

Table~\ref{tab:calibration} demonstrates that SDE-TGNN produces well-calibrated uncertainty estimates. Expected Calibration Error (ECE) of 0.042 indicates predicted confidences closely match empirical accuracy. This is substantially better than MC Dropout (0.092) and Deep Ensemble (0.078), while being computationally more efficient (single forward pass vs. 50 samples or 5 models).

The Fokker-Planck moment propagation provides principled uncertainty quantification grounded in the stochastic dynamics, whereas dropout and ensembles approximate uncertainty indirectly through model variation.

\subsection{Adversarial Robustness}

\begin{table}[t]
\centering
\caption{Adversarial Robustness (Accuracy \% under attack)}
\label{tab:adversarial}
\small
\begin{tabular}{@{}lcccc@{}}
\toprule
\textbf{Model} & \textbf{Clean} & \textbf{FGSM-0.05} & \textbf{PGD-20} & \textbf{Cert-0.25} \\
\midrule
BiLSTM & 93.8 & 67.2 & 52.4 & -- \\
Neural ODE & 96.0 & 71.8 & 58.9 & -- \\
Deep Ensemble & 95.7 & 74.5 & 63.2 & -- \\
\textbf{SDE-TGNN} & \textbf{99.4} & \textbf{89.7} & \textbf{84.3} & \textbf{76.2} \\
\bottomrule
\end{tabular}
\end{table}

Table~\ref{tab:adversarial} shows SDE-TGNN exhibits strong adversarial robustness. Under FGSM attack with $\epsilon=0.05$, accuracy remains 89.7\% compared to 74.5\% for the next-best baseline. Under 20-iteration PGD attack, accuracy is 84.3\% vs. 63.2\% for ensembles.

Most importantly, SDE-TGNN provides \emph{certified} robustness guarantees. At certification radius $R=0.25$, we guarantee 76.2\% of test examples are correctly classified and cannot be attacked with $\ell_2$ perturbations $< 0.25$. Baselines cannot provide such guarantees.

The inherent stochasticity in the SDE formulation acts as implicit adversarial training, while the Fokker-Planck uncertainty quantification enables efficient randomized smoothing certification.

\subsection{Ablation Studies}

\begin{table}[t]
\centering
\caption{Ablation Study on CIC-IDS2018 Dataset}
\label{tab:ablation}
\small
\begin{tabular}{@{}lcc@{}}
\toprule
\textbf{Configuration} & \textbf{F1 (\%)} & \textbf{ECE} \\
\midrule
Full SDE-TGNN & \textbf{99.4} & \textbf{0.042} \\
\quad w/o multi-scale fusion & 97.8 & 0.051 \\
\quad w/o diffusion (ODE only) & 96.2 & 0.089 \\
\quad w/o graph attention & 95.6 & 0.047 \\
\quad w/o Fokker-Planck & 99.3 & 0.068 \\
\quad w/ single scale ($\tau=2$s) & 98.1 & 0.048 \\
\bottomrule
\end{tabular}
\end{table}

Table~\ref{tab:ablation} presents ablation results:

\textbf{Multi-scale fusion:} Removing multi-scale integration drops F1 by 1.6\%, demonstrating the importance of capturing diverse temporal patterns.

\textbf{Stochastic diffusion:} Using deterministic ODE dynamics (removing diffusion) reduces F1 by 3.2\% and significantly increases ECE, confirming that stochastic modeling is essential for both accuracy and calibration.

\textbf{Graph structure:} Removing graph attention drops F1 by 3.8\%, showing that network topology information is valuable.

\textbf{Fokker-Planck propagation:} Replacing analytical moment propagation with Monte Carlo sampling (50 samples) maintains similar F1 but increases ECE by 0.026, validating the efficiency of our approach.

\textbf{Single scale:} Using only medium-term scale reduces F1 by 1.3\%, indicating that both short and long-term patterns contribute to detection.

\subsection{Computational Efficiency}

\begin{table}[t]
\centering
\caption{Computational Efficiency Comparison}
\label{tab:efficiency}
\small
\begin{tabular}{@{}lccc@{}}
\toprule
\textbf{Model} & \textbf{Params} & \textbf{Train (s/epoch)} & \textbf{Infer (ms)} \\
\midrule
BiLSTM & 1.2M & 45 & 2.3 \\
Neural ODE & 2.1M & 187 & 12.7 \\
MC Dropout & 1.2M & 45 & 115 (50 samples) \\
Deep Ensemble & 6.0M & 225 (5 models) & 11.5 (5 models) \\
\textbf{SDE-TGNN} & \textbf{3.8M} & \textbf{156} & \textbf{18.4} \\
\bottomrule
\end{tabular}
\end{table}

Table~\ref{tab:efficiency} compares computational costs. SDE-TGNN requires moderate training time (156s/epoch), between BiLSTM (45s) and Deep Ensemble (225s). Inference latency is 18.4ms, far faster than MC Dropout's 115ms (50 forward passes) and competitive with Neural ODE (12.7ms). The analytical Fokker-Planck propagation avoids the repeated sampling overhead of Monte Carlo methods, making SDE-TGNN practical for real-time deployment.

\subsection{Cross-Domain Generalization}

To evaluate cross-domain transfer, we train on three domains (ICS3D: Cloud + IIoT + Container) and test on the other three (IIS3D: Enterprise + IoT + Hybrid) without fine-tuning:

\begin{table}[t]
\centering
\caption{Cross-Domain Transfer Performance (F1 \%)}
\label{tab:transfer}
\small
\begin{tabular}{@{}lcccc@{}}
\toprule
\textbf{Model} & \textbf{IDS18} & \textbf{IoT23} & \textbf{UNSW} & \textbf{Avg} \\
\midrule
BiLSTM & 78.3 & 74.1 & 76.9 & 76.4 \\
Neural ODE & 82.7 & 79.5 & 81.2 & 81.1 \\
\textbf{SDE-TGNN} & \textbf{91.5} & \textbf{89.8} & \textbf{90.3} & \textbf{90.5} \\
\bottomrule
\end{tabular}
\end{table}

SDE-TGNN achieves 90.5\% average F1 on unseen domains, substantially outperforming baselines (81.1\% for Neural ODE, 76.4\% for BiLSTM). This demonstrates strong cross-domain generalization, attributed to the continuous-time formulation and multi-scale fusion that capture domain-invariant temporal dynamics.

\section{Discussion}

\subsection{Key Findings}

Our experimental results demonstrate three primary advantages of SDE-TGNN:

\textbf{(1) Superior Detection Performance:} The continuous-time stochastic modeling captures both deterministic attack patterns and stochastic traffic variability, achieving 98.7\% weighted F1 across six diverse domains. The multi-scale temporal fusion enables simultaneous detection of attacks with different temporal signatures.

\textbf{(2) Principled Uncertainty Quantification:} Analytical Fokker-Planck moment propagation provides calibrated uncertainty estimates (ECE $<$ 0.05) with $O(d^2)$ complexity, far more efficient than sampling-based methods. This enables risk-aware decision making in production security systems.

\textbf{(3) Certified Robustness:} The inherent stochasticity enables randomized smoothing-based certification, providing provable defense guarantees against $\ell_2$-bounded adversaries. SDE-TGNN maintains 76.2\% certified accuracy at radius 0.25, while baselines cannot provide certificates.

\subsection{Practical Implications}

SDE-TGNN addresses critical requirements for production intrusion detection:

\textbf{Unified Architecture:} A single model handles heterogeneous domains (cloud, IoT, enterprise, etc.) without domain-specific modifications or fine-tuning, reducing deployment complexity.

\textbf{Real-Time Inference:} 18.4ms inference latency with uncertainty quantification enables real-time monitoring of high-throughput networks.

\textbf{Confidence Estimation:} Calibrated uncertainty scores support risk-aware alerting (high-confidence detections trigger immediate responses, low-confidence detections queue for analyst review).

\textbf{Adversarial Defense:} Certified robustness guarantees provide assurance against evasion attacks, critical for safety-critical deployments.

\subsection{Limitations and Future Work}

\textbf{Scalability:} While our moment propagation approach is efficient ($O(d^2)$), scaling to very large graphs (millions of nodes) requires further optimization, potentially through graph sampling or hierarchical representations.

\textbf{Interpretability:} The learned drift and diffusion functions are neural networks, limiting interpretability. Future work could explore symbolic regression or attention visualization techniques to understand learned dynamics.

\textbf{Non-Gaussian Uncertainty:} Our Fokker-Planck propagation assumes Gaussian state distributions via linearization. For highly nonlinear dynamics, higher-order moment propagation or particle filtering may be beneficial.

\textbf{Adaptive Attacks:} While we provide certified robustness against $\ell_2$-bounded perturbations, adaptive attacks targeting the stochastic formulation or temporal structure merit investigation.

\textbf{Online Learning:} The current framework uses batch training. Extending to online/continual learning would enable adaptation to evolving attack patterns without full retraining.

\section{Conclusion}

We introduced SDE-TGNN, a novel deep learning framework for network intrusion detection that combines Temporal Graph Neural Networks with Stochastic Differential Equations and Fokker-Planck-based uncertainty propagation. By modeling network state evolution as continuous-time stochastic dynamics with learned drift and diffusion functions, SDE-TGNN achieves state-of-the-art detection performance (98.7\% weighted F1) across six heterogeneous network security domains totaling 53+ million samples. The analytical moment propagation approach provides calibrated uncertainty estimates with $O(d^2)$ efficiency, orders of magnitude faster than sampling-based methods. The inherent stochasticity enables certified adversarial robustness guarantees via randomized smoothing, addressing critical requirements for safety-critical security deployments.

Our comprehensive evaluation demonstrates that SDE-TGNN outperforms eight baseline models spanning classical ML, deep learning, graph neural networks, continuous-depth models, and Bayesian approximations. Ablation studies confirm the importance of multi-scale temporal fusion, stochastic dynamics, and graph structure. Cross-domain transfer experiments show strong generalization to unseen network domains without fine-tuning.

The unified architecture, real-time inference capability, principled uncertainty quantification, and certified robustness make SDE-TGNN a practical solution for real-world heterogeneous network security infrastructure. We release our implementation, datasets, and evaluation protocols to facilitate reproducibility and future research in stochastic deep learning for cybersecurity.

\section*{Acknowledgments}

This research was supported by the National Research Nuclear University MEPhI. We thank the anonymous reviewers for their constructive feedback.

\bibliographystyle{IEEEtran}
\bibliography{references}

\end{document}
